\documentclass{../../../templates/natsirt}

\usepackage{pgfplots}
\pgfplotsset{compat=1.18}

\title{Linear Algebra 2}
\author{Yi Hong}
\subtitle{Note:p}

\setlength{\parindent}{0pt}
\setlength{\parskip}{3pt}
\setlength{\headheight}{15pt}

% Custom Box Environments for Pure Active-Recall Checklist
\newcounter{topic}[section]

\newenvironment{recallbox}[1][]{%
    \begin{bluebox}[#1]\refstepcounter{topic}\textbf{Definitions \& Concepts to State \thesection.\thetopic}%
}{\end{bluebox}}

\newenvironment{theoremsbox}[1][]{%
    \begin{greenbox}[#1]\refstepcounter{topic}\textbf{Theorems \& Identities to State \thesection.\thetopic}%
}{\end{greenbox}}

\newenvironment{challengebox}[1][]{%
    \begin{redbox}[#1]\refstepcounter{topic}\textbf{Proof Challenges \& Calculations \thesection.\thetopic}%
}{\end{redbox}}

\newenvironment{warningbox}[1][]{\begin{whitebox}[#1]\textbf{Conventions \& Exam Pitfalls. }}{\end{whitebox}}
\newenvironment{reviewbox}[1][]{\begin{whitebox}[#1]\textbf{Textbook Exercise Checklist. }}{\end{whitebox}}

\newcommand{\citem}{\item[$\square$]}
\newcommand{\citemdone}{\item[$\text{\rlap{$\checkmark$}}\square$]}

\DeclareMathOperator{\proj}{proj}
\DeclareMathOperator{\tr}{tr}
\DeclareMathOperator{\Row}{row}
\DeclareMathOperator{\Col}{col}
\DeclareMathOperator{\Null}{null}
\DeclareMathOperator{\rank}{rank}
\DeclareMathOperator{\Span}{span}
\DeclareMathOperator{\diag}{diag}

\begin{document}

% =========================================================================
% CHAPTER 0: JORDAN CANONICAL FORM
% =========================================================================
\setcounter{section}{-1}
\section{Jordan Canonical Form}

\begin{recallbox}
    \begin{itemize}
        \citem \textbf{Jordan Block:} Write the form of a Jordan block $J_k(\lambda)$.
        \citem \textbf{Jordan Canonical Form:} What does it mean for a matrix to be written in Jordan Canonical Form ($A = P J P^{-1}$)?
    \end{itemize}
\end{recallbox}

\begin{warningbox}
    \begin{itemize}
        \citem \textbf{Universality of Jordan Canonical Form:} Notice that every square matrix is similar to a matrix in Jordan canonical form (over $\C$, or any algebraically closed field).
    \end{itemize}
\end{warningbox}

\begin{reviewbox}
    \begin{itemize}
        \citem \textbf{Lecturer Notes:} Back to refer to all the exercises that appear in the lecturer's notes.
        \citem \textbf{Exercise Revision:} Review all exercises for Chapter 0.
    \end{itemize}
\end{reviewbox}


% =========================================================================
% CHAPTER 6: INNER PRODUCT SPACES
% =========================================================================
\setcounter{section}{5}
\section{Inner Product Spaces}

% -------------------------------------------------------------------------
% SECTION 6.1: REAL INNER PRODUCT SPACES
% -------------------------------------------------------------------------
\subsection{Real Inner Product Spaces}

\begin{recallbox}
    \begin{itemize}
        \citem \textbf{Conditions of a Real Inner Product Space:} Write the four conditions (axioms: symmetry, additivity, homogeneity, positivity) of a real inner product space.
        \citem \textbf{Norm and Distance:} Write the definition of norm $\|u\|$ and distance $d(u, v)$.
        \citem \textbf{Weighted Inner Product:} What is the weighted inner product?
    \end{itemize}
\end{recallbox}

\begin{challengebox}
    \begin{itemize}
        \citem \textbf{Properties of Norm and Distance:} Prove that if $u$ and $v$ are vectors in a real inner product space $V$, and if $k$ is a scalar, then:
        \begin{itemize}
            \citem (a) $\|v\| \ge 0$, with equality if and only if $v = 0$.
            \citem (b) $\|kv\| = |k| \|v\|$.
            \citem (c) $d(u, v) = d(v, u)$.
            \citem (d) $d(u, v) \ge 0$, with equality if and only if $u = v$.
        \end{itemize}
        \citem \textbf{Algebraic Properties of Inner Products:} Prove that if $u, v$, and $w$ are vectors in a real inner product space $V$, and if $k$ is a scalar, then:
        \begin{itemize}
            \citem (a) $\langle 0, v \rangle = \langle v, 0 \rangle = 0$.
            \citem (b) $\langle u, v + w \rangle = \langle u, v \rangle + \langle u, w \rangle$.
            \citem (c) $\langle u, v - w \rangle = \langle u, v \rangle - \langle u, w \rangle$.
            \citem (d) $\langle u - v, w \rangle = \langle u, w \rangle - \langle v, w \rangle$.
            \citem (e) $k \langle u, v \rangle = \langle u, kv \rangle$.
        \end{itemize}
        \citem \textbf{Inner Product Expansion:} Show that $\langle u - 2v, 3u + 4v \rangle = 3\|u\|^2 - 2\langle u, v \rangle - 8\|v\|^2$.
    \end{itemize}
\end{challengebox}

\begin{recallbox}
    \begin{itemize}
        \citem \textbf{Inner Product Generated by Matrices:} Write the definition of the inner product generated by matrices ($\langle u, v \rangle = Au \cdot Av = u^T A^T A v$ for invertible $A$), and verify it is indeed an inner product space.
        \citem \textbf{Standard Inner Product on $M_{nn}$:} Write the definition of the standard inner product on $M_{nn}$ ($\langle U, V \rangle = \tr(U^T V) = \sum_{i,j} u_{ij} v_{ij}$), and verify it is indeed an inner product space.
        \citem \textbf{Standard Inner Product on $P_n$:} Write the definition of the standard inner product on $P_n$ ($\langle p, q \rangle = a_0 b_0 + a_1 b_1 + \dots + a_n b_n$), and verify it is indeed an inner product space.
        \citem \textbf{Evaluation Inner Product on $P_n$:} Write the definition of the evaluation inner product on $P_n$ ($\langle p, q \rangle = p(x_0)q(x_0) + \dots + p(x_n)q(x_n)$ at $n+1$ distinct points), and verify it is indeed an inner product space.
        \citem \textbf{Inner Product on $C[a, b]$:} Write the definition of the inner product on $C[a, b]$ ($\langle f, g \rangle = \int_a^b f(x) g(x) \, dx$), and verify it is indeed an inner product space.
    \end{itemize}
\end{recallbox}

\begin{reviewbox}
    \begin{itemize}
        \citem \textbf{Section 6.1 Exercise Revision:}
        \begin{itemize}
            \citem Exercise 3
            \citem Exercise 5
            \citem Exercise 9
            \citem Exercise 21
            \citem Exercise 23
            \citem Exercise 34
        \end{itemize}
    \end{itemize}
\end{reviewbox}

% -------------------------------------------------------------------------
% SECTION 6.2: ANGLE AND ORTHOGONALITY IN INNER PRODUCT SPACES
% -------------------------------------------------------------------------
\subsection{Angle and Orthogonality in Inner Product Spaces}

\begin{theoremsbox}
    \begin{itemize}
        \citem \textbf{Cauchy-Schwarz Inequality:} Prove the Cauchy-Schwarz Inequality: If $u$ and $v$ are vectors in a real inner product space $V$, then
        \[ |\langle u, v \rangle| \le \|u\| \|v\|. \]
    \end{itemize}
\end{theoremsbox}

\begin{recallbox}
    \begin{itemize}
        \citem \textbf{Angle Between Vectors:} For $u$ and $v$ two non-zero vectors in an inner product space, how do we define the angle between them ($\cos\theta = \frac{\langle u, v \rangle}{\|u\| \|v\|}$), and also write the range of angle ($0 \le \theta \le \pi$).
        \citem \textbf{Orthogonality:} Write the definition of when two vectors in an inner product space are called orthogonal ($\langle u, v \rangle = 0$).
        \citem \textbf{Orthogonal Complement:} Suppose $W$ is a subspace of $V$, write the definition of the orthogonal complement of $W$, denoted $W^\perp = \{v \in V \mid \langle v, w \rangle = 0 \text{ for all } w \in W\}$.
    \end{itemize}
\end{recallbox}

\begin{challengebox}
    \begin{itemize}
        \citem \textbf{Triangle Inequalities:} Suppose $u, v, w$ are vectors in an inner product space. Prove that:
        \begin{itemize}
            \citem $\|u + v\| \le \|u\| + \|v\|$
            \citem $d(u, v) \le d(u, w) + d(w, v)$
        \end{itemize}
        \citem \textbf{Generalized Theorem of Pythagoras:} Prove the generalized theorem of Pythagoras: $u$ and $v$ are orthogonal if and only if $\|u + v\|^2 = \|u\|^2 + \|v\|^2$.
        \citem \textbf{Subspace Property of $W^\perp$:} Prove that $W^\perp$ is a subspace of $V$.
        \citem \textbf{Intersection with Orthogonal Complement:} Prove that $W \cap W^\perp = \{0\}$.
        \citem \textbf{Invertibility of the Gram Matrix:} Prove that: Let $W$ be a finite-dimensional subspace of an inner product space $V$, and let $\{w_1, w_2, \dots, w_k\}$ be a basis of $W$. Define the $k \times k$ matrix $A = (a_{ij})$ by
        \[ a_{ij} = \langle w_i, w_j \rangle. \]
        Then $A$ is invertible.
        \citem \textbf{Dimension Formula for Orthogonal Complements:} Prove that: Suppose $V$ is finite-dimensional, and $W$ is a subspace of $V$, then
        \[ \dim W + \dim W^\perp = \dim V. \]
        \citem \textbf{Double Orthogonal Complement:} Show that $(W^\perp)^\perp = W$.
    \end{itemize}
\end{challengebox}

\begin{warningbox}
    \begin{itemize}
        \citem \textbf{Fundamental Subspaces Dualities:} Notice that under the standard inner product:
        \begin{itemize}
            \citem $\Row(A) = W \iff \Null(A) = W^\perp$
            \citem $\Col(A) = W \iff \Null(A^T) = W^\perp$
        \end{itemize}
    \end{itemize}
\end{warningbox}

\begin{reviewbox}
    \begin{itemize}
        \citem \textbf{Section 6.2 Exercise Revision:}
        \begin{itemize}
            \citem Exercise 1
            \citem Exercise 5
            \citem Exercise 9
            \citem Exercise 16
            \citem Exercise 18
            \citem Exercise 29 (see the way)
            \citem Exercise 45
            \citem Exercise 49
        \end{itemize}
    \end{itemize}
\end{reviewbox}

% -------------------------------------------------------------------------
% SECTION 6.3: GRAM-SCHMIDT PROCESS; QR-DECOMPOSITION
% -------------------------------------------------------------------------
\subsection{Gram-Schmidt Process; QR-Decomposition}

\begin{recallbox}
    \begin{itemize}
        \citem \textbf{Orthogonal and Orthonormal Sets:} Write the definition of when a set is orthogonal, and when the set is orthonormal.
        \citem \textbf{Orthogonal Basis and Orthonormal Basis:} Write the definition of an orthogonal basis and an orthonormal basis.
        \citem \textbf{Gram-Schmidt Process:} To convert the basis $\{u_1, \dots, u_r\}$ to an orthogonal basis $\{v_1, \dots, v_r\}$, we apply the Gram-Schmidt process. Try the full step of it:
        \begin{align*}
            v_1 &= u_1, \\
            v_2 &= u_2 - \frac{\langle u_2, v_1 \rangle}{\|v_1\|^2} v_1, \\
            v_3 &= u_3 - \frac{\langle u_3, v_1 \rangle}{\|v_1\|^2} v_1 - \frac{\langle u_3, v_2 \rangle}{\|v_2\|^2} v_2, \\
            &\;\;\vdots \\
            v_r &= u_r - \sum_{j=1}^{r-1} \frac{\langle u_r, v_j \rangle}{\|v_j\|^2} v_j.
        \end{align*}
        And normalize each $v_i$ to obtain an orthonormal basis: $q_i = \frac{v_i}{\|v_i\|}$.
    \end{itemize}
\end{recallbox}

\begin{theoremsbox}
    \begin{itemize}
        \citem \textbf{Linear Independence of Orthogonal Sets:} Prove that if $S = \{v_1, v_2, \dots, v_n\}$ is an orthogonal set of non-zero vectors in an inner product space $V$, then $S$ is linearly independent.
        \citem \textbf{Fourier Expansion Formula:} Prove that if $S = \{v_1, v_2, \dots, v_n\}$ is an orthogonal basis for an inner product space $V$, and if $u$ is any vector in $V$, then
        \[ u = \frac{\langle u, v_1 \rangle}{\langle v_1, v_1 \rangle} v_1 + \dots + \frac{\langle u, v_n \rangle}{\langle v_n, v_n \rangle} v_n. \]
        \citem \textbf{The Projection Theorem:} Write the Projection Theorem, and prove the uniqueness of expression ($u = w_1 + w_2$ with $w_1 \in W$ and $w_2 \in W^\perp$).
        \citem \textbf{Orthogonal Projection Formulas:} Let $W$ be a finite-dimensional subspace of an inner product space $V$, prove that:
        \begin{itemize}
            \citem (a) If $\{v_1, v_2, \dots, v_r\}$ is an orthogonal basis for $W$, and if $u$ is any vector in $V$, then
            \[ \proj_W u = \frac{\langle u, v_1 \rangle}{\langle v_1, v_1 \rangle} v_1 + \dots + \frac{\langle u, v_r \rangle}{\langle v_r, v_r \rangle} v_r \]
            \citem (b) If $\{v_1, v_2, \dots, v_r\}$ is an orthonormal basis for an inner product space $W$, and if $u$ is any vector in $V$, then
            \[ \proj_W u = \langle u, v_1 \rangle v_1 + \dots + \langle u, v_r \rangle v_r \]
        \end{itemize}
        \citem \textbf{Enlarging Orthogonal and Orthonormal Sets:} Prove that if $W$ is a finite-dimensional inner product space, then:
        \begin{itemize}
            \citem (a) Every orthogonal set of non-zero vectors in $W$ can be enlarged to an orthogonal basis for $W$.
            \citem (b) Every orthonormal set in $W$ can be enlarged to an orthonormal basis for $W$.
        \end{itemize}
        \citem \textbf{QR-Decomposition Theorem:} Prove the QR-decomposition: If $A$ is an $m \times n$ matrix with linearly independent column vectors, then $A$ can be factored as
        \[ A = QR \]
        where $Q$ is an $m \times n$ matrix with orthonormal column vectors, and $R$ is an $n \times n$ invertible upper triangular matrix.
    \end{itemize}
\end{theoremsbox}

\begin{reviewbox}
    \begin{itemize}
        \citem \textbf{Section 6.3 Exercise Revision:}
        \begin{itemize}
            \citem Exercise 3
            \citem Exercise 7
            \citem Exercise 19
            \citem Exercise 25
            \citem Exercise 31
            \citem Exercise 35
            \citem Exercise 48
            \citem Exercise 49
            \citem Exercise 55
        \end{itemize}
    \end{itemize}
\end{reviewbox}

% -------------------------------------------------------------------------
% SECTION 6.4: BEST APPROXIMATION; LEAST SQUARES
% -------------------------------------------------------------------------
\subsection{Best Approximation; Least Squares}

\begin{recallbox}
    \begin{itemize}
        \citem \textbf{Least Squares Fundamentals:} Write the definition of the least squares problem, and also the definitions of:
        \begin{itemize}
            \citem Least squares solution $\hat{x}$
            \citem Least squares error vector $b - A\hat{x}$
            \citem Least squares error $\|b - A\hat{x}\|$
        \end{itemize}
    \end{itemize}
\end{recallbox}

\begin{theoremsbox}
    \begin{itemize}
        \citem \textbf{Best Approximation Theorem:} Prove the Best Approximation Theorem: If $W$ is a finite-dimensional subspace of an inner product space $V$, and $u \in V$, then for every $w \in W$ such that $w \ne \proj_W u$, we have
        \[ \|u - \proj_W u\| < \|u - w\|. \]
        \citem \textbf{Normal System Theorem:} Prove the theorem: For every linear system $Ax = b$, the associated normal system
        \[ A^T A x = A^T b \]
        is consistent, and all solutions of this system are least squares solutions of $Ax = b$. Moreover, if $W$ is the column space of $A$, and $x$ is any least squares solution of $Ax = b$, then the orthogonal projection of $b$ on $W$ is
        \[ \proj_W b = Ax. \]
        \citem \textbf{Invertibility Criterion for $A^T A$:} Prove that if $A$ is an $m \times n$ matrix, then the following are equivalent:
        \begin{itemize}
            \citem (a) The column vectors of $A$ are linearly independent.
            \citem (b) $A^T A$ is invertible.
        \end{itemize}
        \citem \textbf{Unique Least Squares Solution and Projection Matrix:} By the previous theorem, show that the unique least squares solution of $Ax = b$ is
        \[ x = (A^T A)^{-1} A^T b \]
        and
        \[ \proj_W b = A (A^T A)^{-1} A^T b. \]
        \citem \textbf{Least Squares via QR-Decomposition:} Prove that if $A$ is an $m \times n$ matrix with linearly independent column vectors, and if $A = QR$ is a QR-decomposition of $A$, then for each $b$ in $\R^m$, the system $Ax = b$ has a unique least squares solution given by
        \[ x = R^{-1} Q^T b. \]
    \end{itemize}
\end{theoremsbox}

\begin{challengebox}
    \begin{itemize}
        \citem \textbf{Projection Matrix onto a Line in $\R^2$:} Find the standard matrix for the orthogonal projection onto the line $W$ through the origin of $\R^2$ that makes an angle $\theta$ with the positive $x$-axis:
        \[ P_\theta = \begin{bmatrix} \cos^2\theta & \sin\theta\cos\theta \\ \sin\theta\cos\theta & \sin^2\theta \end{bmatrix}. \]
    \end{itemize}
\end{challengebox}

\begin{reviewbox}
    \begin{itemize}
        \citem \textbf{Section 6.4 Exercise Revision:}
        \begin{itemize}
            \citem Exercise 3
            \citem Exercise 11
            \citem Exercise 15
            \citem Exercise 21
            \citem Exercise 25
            \citem Exercise 28
        \end{itemize}
    \end{itemize}
\end{reviewbox}

% -------------------------------------------------------------------------
% SECTION 6.5: MATHEMATICAL MODELING VIA LEAST SQUARES
% -------------------------------------------------------------------------
\subsection{Mathematical Modeling via Least Squares}

\begin{recallbox}
    \begin{itemize}
        \citem \textbf{Regression Line Concept:} Write the concept of how we find the regression line ($y = a + bx$) via normal equations.
        \citem \textbf{Higher Degree Polynomial Fitting:} What if we do not want to use a regression line to fit the data points, but instead use some higher degree polynomial? Write the procedure to find the polynomial.
    \end{itemize}
\end{recallbox}

\begin{challengebox}
    \begin{itemize}
        \citem \textbf{Straight Line Fitting Calculation:} Find the least squares straight line fit to the four points $(0, 1)$, $(1, 3)$, $(2, 4)$, and $(3, 4)$.
    \end{itemize}
\end{challengebox}

\begin{reviewbox}
    \begin{itemize}
        \citem \textbf{Section 6.5 Exercise Revision:}
        \begin{itemize}
            \citem Exercise 3
        \end{itemize}
    \end{itemize}
\end{reviewbox}


% =========================================================================
% CHAPTER 7: DIAGONALIZATION
% =========================================================================
\setcounter{section}{6}
\section{Diagonalization}

% -------------------------------------------------------------------------
% SECTION 7.1: ORTHOGONAL MATRICES
% -------------------------------------------------------------------------
\subsection{Orthogonal Matrices}

\begin{recallbox}
    \begin{itemize}
        \citem \textbf{Orthogonal Matrix Definition:} Write the definition of an orthogonal matrix ($A^{-1} = A^T$, or $A^T A = A A^T = I$).
        \citem \textbf{3D Rotation Matrices:} Write the rotation matrices for 3D (rotating a vector with respect to the $x$-, $y$-, and $z$-axes):
        \begin{align*}
            R_x(\theta) &= \begin{bmatrix} 1 & 0 & 0 \\ 0 & \cos\theta & -\sin\theta \\ 0 & \sin\theta & \cos\theta \end{bmatrix}, \\[4pt]
            R_y(\theta) &= \begin{bmatrix} \cos\theta & 0 & \sin\theta \\ 0 & 1 & 0 \\ -\sin\theta & 0 & \cos\theta \end{bmatrix}, \\[4pt]
            R_z(\theta) &= \begin{bmatrix} \cos\theta & -\sin\theta & 0 \\ \sin\theta & \cos\theta & 0 \\ 0 & 0 & 1 \end{bmatrix}.
        \end{align*}
    \end{itemize}
\end{recallbox}

\begin{theoremsbox}
    \begin{itemize}
        \citem \textbf{Characterizations of Orthogonal Matrices:} Prove that the following are equivalent for an $n \times n$ matrix $A$:
        \begin{itemize}
            \citem (a) $A$ is orthogonal.
            \citem (b) The row vectors of $A$ form an orthonormal set in $\R^n$ with the Euclidean inner product.
            \citem (c) The column vectors of $A$ form an orthonormal set in $\R^n$ with the Euclidean inner product.
        \end{itemize}
        \citem \textbf{Four Fundamental Theorems of Orthogonal Matrices:} Prove these four theorems:
        \begin{itemize}
            \citem (a) The transpose of an orthogonal matrix is orthogonal.
            \citem (b) The inverse of an orthogonal matrix is orthogonal.
            \citem (c) A product of orthogonal matrices is orthogonal.
            \citem (d) If $A$ is orthogonal, then $\det(A) = \pm 1$.
        \end{itemize}
        \citem \textbf{Geometric Equivalences (Isometry):} Prove that if $A$ is an $n \times n$ matrix, then the following are equivalent:
        \begin{itemize}
            \citem (a) $A$ is orthogonal.
            \citem (b) $\|Ax\| = \|x\|$ for all $x \in \R^n$.
            \citem (c) $Ax \cdot Ay = x \cdot y$ for all $x$ and $y$ in $\R^n$.
        \end{itemize}
        \citem \textbf{Coordinate Properties in Orthonormal Bases:} Prove that if $S$ is an orthonormal basis for an $n$-dimensional inner product space $V$, and if
        \[ (u)_S = (u_1, u_2, \dots, u_n) \quad \text{and} \quad (v)_S = (v_1, v_2, \dots, v_n), \]
        then:
        \begin{itemize}
            \citem (a) $\|u\| = \sqrt{u_1^2 + u_2^2 + \dots + u_n^2}$
            \citem (b) $d(u, v) = \sqrt{(u_1 - v_1)^2 + (u_2 - v_2)^2 + \dots + (u_n - v_n)^2}$
            \citem (c) $\langle u, v \rangle = u_1 v_1 + u_2 v_2 + \dots + u_n v_n$
        \end{itemize}
        \citem \textbf{Transition Matrix Between Orthonormal Bases:} Let $V$ be a finite-dimensional inner product space. Prove that if $P$ is the transition matrix from one orthonormal basis for $V$ to another orthonormal basis for $V$, then $P$ is an orthogonal matrix.
    \end{itemize}
\end{theoremsbox}

\begin{reviewbox}
    \begin{itemize}
        \citem \textbf{Section 7.1 Exercise Revision:} Practice and review all assigned textbook exercises for Section 7.1.
    \end{itemize}
\end{reviewbox}

% -------------------------------------------------------------------------
% SECTION 7.2: ORTHOGONAL DIAGONALIZATION
% -------------------------------------------------------------------------
\subsection{Orthogonal Diagonalization}

\begin{recallbox}
    \begin{itemize}
        \citem \textbf{Orthogonally Similar Matrices:} Write the definition of when $A$ and $B$ are orthogonally similar matrices ($B = P^T A P = P^{-1} A P$ for some orthogonal matrix $P$).
        \citem \textbf{Procedure of Orthogonally Diagonalizing a Symmetric Matrix:} Write the procedure of orthogonally diagonalizing a symmetric matrix:
        \begin{itemize}
            \citem Find all eigenvalues of $A$ by solving $\det(\lambda I - A) = 0$.
            \citem Find a basis for each eigenspace by solving $(\lambda I - A)x = 0$.
            \citem Apply the Gram-Schmidt process to each eigenspace basis with dimension $\ge 2$ to obtain an orthogonal basis.
            \citem Normalize each eigenvector to obtain an orthonormal set of $n$ eigenvectors $u_1, \dots, u_n$.
            \citem Form the orthogonal matrix $P = [u_1 \mid u_2 \mid \dots \mid u_n]$ that diagonalizes $A$, so that $P^T A P = D$.
        \end{itemize}
        \citem \textbf{Spectral Decomposition of a Symmetric Matrix:} Write the spectral decomposition of a symmetric matrix:
        \[ A = \lambda_1 u_1 u_1^T + \lambda_2 u_2 u_2^T + \dots + \lambda_n u_n u_n^T = \sum_{i=1}^n \lambda_i P_i. \]
    \end{itemize}
\end{recallbox}

\begin{theoremsbox}
    \begin{itemize}
        \citem \textbf{Fundamental Theorem of Real Symmetric Matrices:} Prove that if $A$ is an $n \times n$ matrix with real entries, then the following are equivalent:
        \begin{itemize}
            \citem (a) $A$ is orthogonally diagonalizable.
            \citem (b) $A$ has an orthonormal set of $n$ eigenvectors.
            \citem (c) $A$ is symmetric.
        \end{itemize}
        \citem \textbf{Eigenproperties of Symmetric Matrices:} Prove that if $A$ is a symmetric matrix with real entries, then:
        \begin{itemize}
            \citem (a) The eigenvalues of $A$ are all real numbers.
            \citem (b) Eigenvectors from different eigenspaces are orthogonal.
        \end{itemize}
    \end{itemize}
\end{theoremsbox}

\begin{challengebox}
    \begin{itemize}
        \citem \textbf{Diagonalization Calculation:} Find an orthogonal matrix $P$ that diagonalizes
        \[ A = \begin{bmatrix} 4 & 2 & 2 \\ 2 & 4 & 2 \\ 2 & 2 & 4 \end{bmatrix}. \]
    \end{itemize}
\end{challengebox}

\begin{reviewbox}
    \begin{itemize}
        \citem \textbf{Section 7.2 Exercise Revision:} Practice and review all assigned textbook exercises for Section 7.2.
    \end{itemize}
\end{reviewbox}

\end{document}
